% BANKS-SOURCE: pdf=78 print=68 kind=section id=6.3
\subsection{Interaction with heavy fermions: particle paths and external fields}

% BANKS-SOURCE: pdf=78 print=68 kind=prose
If the muon mass were extremely heavy, we would not be able to create muons in low-
energy processes involving electrons, photons, and muons. Furthermore, in this limit,
muons move without much disturbance from their interactions (momentum transfer
much less than the muon mass). Thus we can control the trajectories of any muons and
anti-muons in the initial state, and restrict our attention to trajectories where no anni-
hilation processes are possible. In this section we will discuss a set of approximations
that are useful for describing this limit.

In the path-integral formalism for QED, we can do the integral over the muon field
exactly, with the other fields held fixed, in terms of the Green function of the muon Dirac
operator in an external photon field. If we are trying to evaluate a Green function with
$2N$ muon fields, appropriate for situations in which there are $N$ muons or anti-muons
in the initial state of a scattering process, then the result is

% BANKS-SOURCE: pdf=78 print=68 kind=equation id=6.13
\begin{equation}
  \det\!\left[S(x_i,y_j)\right]\operatorname{Det}\!\left[-\ii\slashed\partial-e\slashed A-M_\mu\right].
  \label{eq:6.13}
\end{equation}
$S(x,y)$ is the solution of
\begin{equation*}
  \left[-\ii\slashed\partial-e\slashed A(x)-M_\mu\right]S(x,y)
  =\ii\delta^D(x-y),
\end{equation*}
with Feynman boundary conditions. The lower-case det refers to the ordinary determi-
nant of the $N\times N$ matrix of propagators, corresponding to all possible contractions of

% BANKS-SOURCE: pdf=79 print=69 kind=prose
initial and final points. The upper-case Det refers to the functional determinant of the
Dirac operator. The functional determinant is negligible as long as all external fields
and momenta are small compared with $M_\mu$.

To evaluate the propagator $S$ we write

% BANKS-SOURCE: pdf=79 print=69 kind=equation id=6.14
\begin{equation}
  S=(-\ii\slashed D-M_\mu)^{-1}
  =(-\ii\slashed D+M_\mu)
  \left(D^2-M_\mu^2+\frac{e}{2}\sigma^{\mu\nu}F_{\mu\nu}+\ii\epsilon\right)^{-1}.
  \label{eq:6.14}
\end{equation}
where $D_\mu=\partial_\mu-\ii e A_\mu$ is the covariant derivative and
$\sigma^{\mu\nu}=\frac{\ii}{2}[\gamma^\mu,\gamma^\nu]$. To evaluate the inverse of the second-order opera-
tor, we think of it as the Hamiltonian for a quantum particle with four position coordi-
nates\footnote[4]{This actually looks a little more conventional in Euclidean space, but we will not bother to go through the analytic continuations.} and a spin, and write the inverse in terms of a mixture of path-integral formalism
for the positional mechanics and a time-ordered formula for the spin dynamics:

% BANKS-SOURCE: pdf=79 print=69 kind=equation id=6.15
\begin{equation}
  \begin{aligned}
  S={}&(-\ii\slashed D+M_\mu)\int_0^\infty ds\,\ee^{-\ii sM_\mu^2}[dx^\mu(s)]
  \exp\!\left\{\ii\int_0^s d\tau\left[\left(\frac{dx}{2d\tau}\right)^2
  +eA_\mu\frac{dx^\mu}{d\tau}\right]\right\}\\[-2pt]
  &\hspace{3cm}\times\mathcal{T}\exp\!\left\{\frac{\ii e}{2}\int_0^s d\tau\,
  \sigma^{\mu\nu}F_{\mu\nu}(x(\tau))\right\}.
  \end{aligned}
  \label{eq:6.15}
\end{equation}
The path integral is over paths satisfying $x(0)=x$, $x(s)=y$. On introducing the
dimensionless variables $sM_\mu^2\equiv u$, $\tau M_\mu^2=v$, it is easy to see that, for large
$M_\mu$, the particle action is large, and is dominated by the free-particle kinetic term, which is of
order $M_\mu^2$. The term involving the line integral of the vector potential is of order one,
while the interaction with the spin is down by $1/M_\mu^2$.

\BanksImplicitHook{I-CH06-004}

Thus we can saturate the path integral by the straight-line path between $x$ and $y$.
The result, to leading order in $M_\mu$, is

% BANKS-SOURCE: pdf=79 print=69 kind=equation id=6.16
\begin{equation}
  S(x,y)=\int_0^\infty ds\,
  \ee^{-\ii sM_\mu^2+\ii(x-y)^2/(4s)}
  \ee^{\frac{\ii e}{s}\int_0^s d\tau\,A_L\!\left(x+\frac{\tau}{s}(y-x)\right)}.
  \label{eq:6.16}
\end{equation}
$A_L\equiv(x-y)^\mu A_\mu$, and the last exponential is just the line integral of the vector
potential along the straight-line particle path and, in particular, it is independent of
the parameter $s$. Thus, the effect of integrating out the massive fermion is to add a
source term to the $A_\mu$ Lagrangian. If we shift the $A_\mu$ field to eliminate the source, it is
shifted exactly by the classical (massive) electromagnetic field produced by the current
of the heavy point particle. Obviously we can repeat this procedure for all the fermion
propagators. To leading order in the mass, the effect of a collection of massive particles
on the rest of the theory is simply to shift the background value of $A_\mu$ from zero to some
other classical field. This justifies the study of such classical external-field problems.
